%% FinWatch Literature Review
%% Consumer Financial Vulnerability Detection: A Machine Learning Perspective
%% Author: Barbara Werobaobayi
%% University of Strathclyde, MSc Data Science
%%
%% Compile on Overleaf: New Project → Upload → literature_review.tex
%% Compiler: pdfLaTeX  |  TeX Live 2023

\documentclass[12pt, a4paper]{article}

%% -- Packages ------------------------------------------------------------------
\usepackage[margin=2.5cm]{geometry}
\usepackage[T1]{fontenc}
\usepackage[utf8]{inputenc}
\usepackage{lmodern}               % Clean vector fonts
\usepackage{microtype}             % Improved justification
\usepackage{setspace}
\usepackage{parskip}               % Space between paragraphs, no indent
\usepackage{titlesec}              % Section heading control
\usepackage{booktabs}              % Professional tables
\usepackage{array}
\usepackage{longtable}
\usepackage{hyperref}
\usepackage[numbers, sort&compress]{natbib}
\usepackage{enumitem}
\usepackage{abstract}
\usepackage{fancyhdr}
\usepackage{amsmath}

%% -- Hyperref (black, no coloured boxes) ------------------------------------─
\hypersetup{
    colorlinks = true,
    linkcolor  = black,
    citecolor  = black,
    urlcolor   = black,
    pdftitle   = {Consumer Financial Vulnerability Detection: A Machine Learning Perspective},
    pdfauthor  = {Barbara Werobaobayi},
}

%% -- Typography --------------------------------------------------------------─
\onehalfspacing
\setlength{\parindent}{0pt}
\setlength{\parskip}{8pt}

\titleformat{\section}{\large\bfseries}{{\thesection.}}{0.6em}{}
\titleformat{\subsection}{\normalsize\bfseries}{{\thesubsection}}{0.5em}{}
\titlespacing{\section}{0pt}{18pt}{6pt}
\titlespacing{\subsection}{0pt}{12pt}{4pt}

%% -- Header / footer ----------------------------------------------------------
\pagestyle{fancy}
\fancyhf{}
\renewcommand{\headrulewidth}{0.4pt}
\fancyhead[L]{\small\textit{FinWatch --- Literature Review}}
\fancyhead[R]{\small\textit{B. Werobaobayi}}
\fancyfoot[C]{\thepage}

%% -- Document ----------------------------------------------------------------─
\begin{document}

\begin{titlepage}
\centering
\vspace*{2cm}

{\LARGE\bfseries Consumer Financial Vulnerability Detection:\\[6pt]
A Machine Learning Perspective\par}

\vspace{1.5cm}

{\large Barbara Werobaobayi\par}
{\normalsize MSc Machine Learning and Deep Learning, University of Strathclyde\par}

\vspace{0.8cm}

{\normalsize \today\par}

\vspace{2cm}

\begin{abstract}
\noindent
Consumer financial vulnerability is an underserved problem in applied machine learning.
Existing credit risk systems are designed to screen new applicants at origination,
leaving institutions with no systematic mechanism for identifying customers
who become vulnerable \emph{after} account opening --- a gap that acquired
regulatory urgency with the introduction of the UK Financial Conduct Authority's
Consumer Duty framework in July 2023.
This review surveys the academic and practitioner literature across six intersecting
domains: the conceptual framing of vulnerability as a predictive target; gradient
boosting methods as the production standard for tabular financial data;
hyperparameter optimisation with Bayesian search; class imbalance treatment;
model explainability and fairness auditing under regulatory constraint;
and the role of macroeconomic context in real-time risk scoring.
We conclude by identifying a persistent gap between regulatory obligation and deployed
capability, and situate \emph{FinWatch} --- a continuous vulnerability scoring
and intervention routing system --- as a direct response to that gap.
\end{abstract}

\vspace{1cm}

\noindent\textbf{Keywords:} financial vulnerability, consumer duty, gradient boosting,
disparate impact, SHAP explainability, drift detection, intervention routing

\vfill
\end{titlepage}

\tableofcontents
\newpage

%% ----------------------------------------------------------------------------─
\section{Introduction}
%% ----------------------------------------------------------------------------─

The dominant paradigm in financial machine learning remains credit default prediction
at origination: given an application, will this customer repay?
This framing is well-studied, commercially mature, and increasingly insufficient.
It answers the wrong question for a large class of modern regulatory obligations.

The Financial Conduct Authority (FCA) Consumer Duty, which took full legal effect
in July 2023, requires UK retail banks to act to ``deliver good outcomes for
retail customers'' and to make ``specific provision for the needs of customers in
vulnerable circumstances'' \cite{fca2023duty}.
The obligation is \emph{prospective and continuous} --- banks must identify customers
who are currently experiencing financial stress, not merely those who presented
as poor credit risks at origination.

This creates demand for a fundamentally different class of model: one that scores
existing customers on a recurring basis, uses features reflecting current economic
conditions, and routes customers to proportionate interventions rather than
binary approve/decline outputs.
This review surveys the literature relevant to building such a system.

The scope covers: gradient boosting methods for tabular data; Bayesian hyperparameter
optimisation; treatment of class imbalance; model calibration; explainability under
Consumer Duty; fairness auditing under the 4/5ths rule; feature drift detection;
and the integration of macroeconomic signals as dynamic covariates.

%% ----------------------------------------------------------------------------─
\section{From Default Prediction to Vulnerability Detection}
%% ----------------------------------------------------------------------------─

\subsection{The Conceptual Shift}

Thakor \cite{thakor2020} traces the theoretical basis of credit scoring from
logistic regression-era probability-of-default models to the modern view that
credit risk is a continuous, time-varying property of a customer's circumstances.
The important practical implication is that customers who were creditworthy at
origination may become vulnerable months or years later --- due to job loss,
relationship breakdown, or macroeconomic shock --- without any new credit event
triggering a re-evaluation.

Kaltwasser and Volpe \cite{kaltwasser2018} distinguish between \emph{structural}
vulnerability (characteristics of the individual, such as age, health, or low
financial literacy) and \emph{situational} vulnerability (transient circumstances
such as bereavement, sudden income loss, or debt accumulation).
Both types are relevant; only situational vulnerability is detectable in real time
from transaction and bureau data.

\subsection{The Label Lag Problem}

A fundamental challenge in building vulnerability models is label availability.
Unlike fraud detection, where ground truth (a confirmed fraudulent transaction) is
available within days, financial distress labels arrive with a lag of 6--18 months
between the onset of stress and a formal credit event such as a missed payment or
default \cite{lessmann2015}.
This means:

\begin{itemize}[itemsep=2pt]
    \item Training sets necessarily reflect historical rather than current conditions
    \item Concept drift --- the change in the statistical relationship between features
          and the target --- is an inherent operational risk rather than an edge case
    \item Threshold calibration must account for the temporal gap between scoring
          and the eventual outcome
\end{itemize}

Brzezinski \cite{brzezinski2014} provides a taxonomy of concept drift in financial
models and demonstrates that periodic retraining is insufficient; drift detection
metrics are needed to trigger unscheduled retraining when the feature distribution
departs materially from the training distribution.

%% ----------------------------------------------------------------------------─
\section{Regulatory Context: FCA Consumer Duty}
%% ----------------------------------------------------------------------------─

\subsection{Legislative Background}

The FCA Consumer Duty (PS22/9) represents the most significant reform of UK conduct
regulation since the Retail Distribution Review \cite{fca2022ps229}.
The Duty imposes four outcome requirements on firms: products and services;
price and value; consumer understanding; and consumer support.
The fourth outcome is directly implicated by vulnerability detection: firms must
provide ``support that meets the needs of customers in vulnerable circumstances''.

Critically, the Duty uses the word ``proactively'' twelve times in the final guidance.
Reactive identification --- waiting for a customer to contact the firm in distress ---
is explicitly insufficient \cite{fca2023vulnerable}.

\subsection{FCA's Definition of Vulnerability}

The FCA defines a vulnerable customer as ``someone who, due to their personal
circumstances, is especially susceptible to harm, particularly when a firm is not
acting with appropriate levels of care'' \cite{fg2120}.
The guidance identifies four drivers of vulnerability: health (physical or mental),
life events, resilience (financial or emotional), and capability (literacy or
digital skills).

This multidimensional definition creates a modelling challenge: no single proxy
variable captures all four drivers, and the most important drivers (health, mental
state) are either unavailable or protected under data protection law.
Practitioners therefore work with observable proxies from bureau data and
transaction history, which capture resilience and some life-event signals
but not health or capability directly.

\subsection{Implications for Model Design}

The regulatory context implies several non-negotiable design constraints:

\begin{enumerate}[itemsep=2pt]
    \item Models must produce \emph{calibrated probabilities}, not just rankings,
          to support risk-proportionate intervention routing
    \item Decisions must be \emph{explainable} to the customer on request
          (FCA and UK GDPR Article 22)
    \item Protected characteristics (gender, age, ethnicity) must not drive
          disproportionate escalation rates
    \item The model must be subject to ongoing \emph{monitoring} with documented
          retrain triggers
\end{enumerate}

%% ----------------------------------------------------------------------------─
\section{Gradient Boosting as the Production Standard}
%% ----------------------------------------------------------------------------─

\subsection{Empirical Superiority on Tabular Data}

The choice of model family for structured financial data is well-settled in the
practitioner literature.
Grinsztajn, Oyallon, and Varoquaux \cite{grinsztajn2022} conduct a systematic
comparison of tree-based methods against deep learning architectures on 45 tabular
datasets and find that tree ensembles --- specifically gradient boosted trees ---
consistently outperform neural networks when $n < 10^6$ and feature engineering
is meaningful.
The result is robust across dataset sizes and preprocessing strategies.

Chen and Guestrin's XGBoost \cite{chenguest2016} and Ke et al.'s LightGBM \cite{ke2017}
are the two dominant implementations.
LightGBM's leaf-wise growth strategy and histogram-based binning make it
materially faster than XGBoost on high-cardinality features common in bureau data,
while XGBoost's column subsampling tends to produce better-calibrated raw probabilities
\cite{borisov2022}.
Running both and selecting a champion on a held-out validation set is therefore
best practice rather than redundancy.

\subsection{Calibration}

Raw gradient boosting scores are not probabilities --- they are log-odds values
that require post-hoc calibration.
Platt scaling (logistic calibration of the raw score) \cite{platt1999} is the
standard approach for credit models because it preserves rank ordering while
correcting the score magnitude.
Niculescu-Mizil and Caruana \cite{niculescu2005} show that Platt scaling is
appropriate for boosting methods specifically, where the raw distribution tends
to be bimodal and away from the probability simplex.
Scikit-learn's \texttt{CalibratedClassifierCV} with \texttt{method='sigmoid'} is
the standard implementation used in production financial ML \cite{pedregosa2011}.

%% ----------------------------------------------------------------------------─
\section{Hyperparameter Optimisation}
%% ----------------------------------------------------------------------------─

\subsection{Bayesian Search over Grid Search}

Grid search and random search were the dominant HPO strategies until 2012.
Bergstra and Bengio \cite{bergstra2012} demonstrated that random search outperforms
grid search in almost all settings because most hyperparameter dimensions contribute
little to performance variation; exploring the important dimensions more densely
is therefore more efficient.

Snoek, Larochelle, and Adams \cite{snoek2012} formalised Bayesian optimisation for
HPO, showing that a Gaussian process surrogate model of the objective landscape
enables directed exploration that outperforms random search, especially when
function evaluation is expensive.

\subsection{Tree-structured Parzen Estimators}

Bergstra, Bardenet, and Bengio \cite{bergstra2011} introduced the Tree-structured
Parzen Estimator (TPE), which models $p(x|y)$ and $p(y)$ separately rather than
fitting a joint Gaussian process.
TPE scales better than GP-based Bayesian optimisation in high-dimensional search
spaces and is the algorithm underlying Optuna \cite{akiba2019}.

Optuna adds two important production features: asynchronous parallel trials
(relevant when training LightGBM on multi-core hardware) and \emph{pruning},
whereby unpromising trials are terminated early based on intermediate validation
metrics, reducing total compute by up to 60\% \cite{akiba2019}.

%% ----------------------------------------------------------------------------─
\section{Class Imbalance}
%% ----------------------------------------------------------------------------─

\subsection{Prevalence in Financial Distress Data}

Financial distress datasets are structurally imbalanced.
In the Home Credit Default Risk dataset \cite{homecredit2018}, the minority class
(default) represents approximately 8\% of observations.
Vulnerability, defined more broadly, may reach 15--20\% but remains a minority
class.
Standard classifiers trained on imbalanced data learn to predict the majority class
and achieve high accuracy while missing nearly all minority-class events ---
the ``accuracy paradox'' \cite{he2009}.

\subsection{SMOTE and Its Constraints}

The Synthetic Minority Over-sampling Technique (SMOTE) \cite{chawla2002} generates
synthetic minority examples by interpolating between existing minority instances
in feature space.
It is the most widely deployed oversampling method in financial ML.
However, two constraints are critical in practice:

First, SMOTE must be applied \emph{only to the training fold} and never to validation
or test data.
Applying SMOTE before the train/test split constitutes a form of data leakage
that inflates reported performance metrics \cite{kapelner2021}.

Second, SMOTE can generate synthetic examples in feature-space regions that are
inconsistent with domain constraints (e.g. negative income values in interpolated
records).
Post-processing validation is therefore recommended \cite{he2009}.

\subsection{Alternative Strategies}

Weiss, McCarthy, and Zabar \cite{weiss2007} survey the alternatives: cost-sensitive
learning (penalising false negatives more heavily), threshold moving (shifting the
decision boundary after training), and ensemble methods (BalancedBaggingClassifier,
EasyEnsemble).
In the vulnerability detection context, cost-sensitive threshold calibration is
preferable to SMOTE for the final decision boundary because it directly encodes
the asymmetric business cost of missing a vulnerable customer versus incorrectly
flagging a non-vulnerable one.

%% ----------------------------------------------------------------------------─
\section{Model Explainability}
%% ----------------------------------------------------------------------------─

\subsection{Regulatory Drivers}

UK GDPR Article 22 grants data subjects the right not to be subject to solely
automated decisions that significantly affect them, and the right to meaningful
information about the logic involved.
For a vulnerability model that determines whether a customer receives proactive
outreach, this right is directly engaged.
The FCA's own guidance on algorithmic decision-making requires that firms be able
to explain to customers why they have been contacted or not contacted as a result
of an automated system \cite{fca2023ai}.

\subsection{SHAP}

Lundberg and Lee \cite{lundberg2017} introduced SHapley Additive exPlanations (SHAP),
grounding feature attribution in cooperative game theory.
SHAP values satisfy three axioms (efficiency, symmetry, and dummy) that no other
attribution method simultaneously satisfies, making them the principled choice
for regulatory environments where the attribution must be defensible.

TreeSHAP \cite{lundberg2020} extends the method to tree ensembles with
$O(TLD^2)$ complexity rather than the exponential cost of exact Shapley computation,
making it practical for production inference at scale.
The combination of LightGBM and TreeSHAP is therefore the natural pairing for
Consumer Duty compliance: the model and its explanation are computed in the same
framework without approximation.

\subsection{Local vs Global Explanations}

Ribeiro, Singh, and Guestrin \cite{ribeiro2016} distinguish between local explanations
(why was this individual customer scored as vulnerable?) and global explanations
(what features drive vulnerability across the portfolio?).
Consumer Duty requires \emph{local} explanations --- the firm must be able to
explain the decision to the specific customer --- while model governance requires
global explanations for validation.
SHAP provides both through the same unified framework.

%% ----------------------------------------------------------------------------─
\section{Fairness and Disparate Impact}
%% ----------------------------------------------------------------------------─

\subsection{The 4/5ths Rule}

The 4/5ths rule (80\% rule) for disparate impact originated in US Equal Employment
Opportunity Commission guidelines \cite{eeoc1978} and has been adopted by the FCA
as an informal threshold for acceptable disparity in automated decision-making
\cite{fca2023ai}.
The Disparate Impact Ratio (DIR) is:

\begin{equation}
    \text{DIR} = \frac{\text{positive outcome rate (disadvantaged group)}}
                      {\text{positive outcome rate (advantaged group)}}
\end{equation}

A DIR below 0.80 indicates that the disadvantaged group receives the positive
outcome (here: escalation to proactive support) at less than 80\% of the rate
of the advantaged group, which constitutes prima facie disparate impact.

\subsection{Sources of Disparity in Financial Models}

Mehrabi et al.\ \cite{mehrabi2021} survey the causes of algorithmic unfairness in
financial contexts.
In vulnerability detection, the primary risk is \emph{historical bias}: if
historically marginalised groups were less likely to receive financial support
(regardless of actual need), a model trained on outcomes will replicate this
pattern.
A secondary risk is \emph{proxy discrimination}: features such as postcode or
employment sector that are not protected attributes but correlate with protected
characteristics.

\subsection{Mitigation Approaches}

Bird et al.'s Fairlearn \cite{bird2020} and IBM's AIF360 \cite{bellamy2019} provide
Python toolkits for both measuring and mitigating disparate impact.
Mitigation strategies fall into three categories: pre-processing (resampling or
re-weighting training data), in-processing (adding fairness constraints to the
objective), and post-processing (adjusting thresholds per group).
For regulated financial models, post-processing threshold adjustment is the most
auditable approach because it leaves the model itself unchanged and the adjustment
is fully documented \cite{hardt2016}.

%% ----------------------------------------------------------------------------─
\section{Macroeconomic Features and Live Data Integration}
%% ----------------------------------------------------------------------------─

\subsection{Why Static Models Fail in Macro Shocks}

A credit model trained in 2019 experienced conditions --- low interest rates,
stable employment, low energy costs --- that have not held since 2022.
The UK interest rate rose from 0.1\% in December 2021 to 5.25\% by August 2023;
CPI peaked at 11.1\% in October 2022 \cite{ons2023}.
A model with no macroeconomic features cannot distinguish a customer who is
financially stressed because of their individual circumstances from one who is
stressed because the entire economy shifted.

This ``macro blindness'' problem means that static models systematically
misclassify customers who were not vulnerable at training time but became
vulnerable due to rate rises --- exactly the population that Consumer Duty
requires proactive identification of.

\subsection{ONS API as a Live Data Source}

The UK Office for National Statistics publishes macroeconomic series via a
public REST API \cite{ons2024api}.
Series relevant to consumer vulnerability include: the Consumer Price Index
(CPIH), unemployment rate (LFS), Bank of England base rate (IUMABEDR), and
household energy expenditure (HHFCE).
These series are updated monthly and can be appended as constant-valued features
to every row in a scoring batch, providing the model with a snapshot of the
current economic environment at the time of scoring.

Butaru et al.\ \cite{butaru2016} demonstrate in a multi-bank US study that macro
variables improve predictive accuracy for financial distress detection by 12--18\%
over bureau-data-only models, with the largest gains during periods of rapid
macroeconomic change.

\subsection{Macro Drift as a Retrain Trigger}

Gama et al.\ \cite{gama2014} formalise the concept of drift detection and propose
that distributional change in input features --- detectable before label drift,
since labels arrive with a lag --- should trigger model re-evaluation.
Storing the training-time distribution of macroeconomic features and computing
z-scores of current values against that baseline provides an early warning system
that is independent of model performance metrics.

%% ----------------------------------------------------------------------------─
\section{Monitoring: Population Stability Index}
%% ----------------------------------------------------------------------------─

The Population Stability Index (PSI) \cite{yurdakul2018} measures distributional
shift between a reference (training) and current (production) dataset:

\begin{equation}
    \text{PSI} = \sum_{i=1}^{k} \left( A_i - E_i \right) \cdot \ln\!\left(\frac{A_i}{E_i}\right)
\end{equation}

where $A_i$ and $E_i$ are the actual and expected proportions in bucket $i$.
The conventional thresholds are:

\begin{center}
\begin{tabular}{ll}
\toprule
PSI & Interpretation \\
\midrule
$< 0.10$ & No significant change \\
$0.10$--$0.20$ & Some change --- monitor \\
$> 0.20$ & Significant change --- investigate \\
$> 0.25$ & Major shift --- retrain \\
\bottomrule
\end{tabular}
\end{center}

PSI is the industry-standard drift metric for financial models because it is
interpretable, requires no labels, and is computed efficiently on the marginal
distribution of each feature independently.
Thomas \cite{thomas2009} discusses PSI in the context of credit scorecard monitoring
and confirms its superiority over KL divergence for one-sided governance reporting.

%% ----------------------------------------------------------------------------─
\section{Frontier: Foundation Models and LLMs in Financial ML}
%% ----------------------------------------------------------------------------─

The most recent wave of research explores whether large language models and
transformer architectures can displace gradient boosting for tabular financial data.

Yu et al.\ \cite{yu2023} introduce FinPT, a framework that serialises tabular
financial records into natural language strings and fine-tunes a language model
for risk prediction.
They report improvements over LightGBM on several benchmark datasets, though the
gains are modest and inference latency is 20--50× higher.

Wang et al.\ \cite{wang2023} propose nuFormer, a transformer architecture with a
dedicated numerical attention head designed to handle the mixed discrete/continuous
nature of credit bureau data without text serialisation.
Initial results are promising on high-cardinality datasets.

The TransactionGPT line of work (Cheng et al.\ \cite{cheng2024}) explores
pre-training on raw transaction sequences, exploiting the natural sequential
structure of spending history to learn latent representations of financial
behaviour without feature engineering.
This approach is compelling for vulnerability detection because spending pattern
changes (grocery substitution towards budget brands, cancellation of discretionary
subscriptions) are early-warning signals that structured features miss.

However, Grinsztajn et al.\ \cite{grinsztajn2022} find that gradient boosting
remains superior on datasets with fewer than $10^6$ observations, which describes
the majority of retail bank vulnerability scoring contexts.
The production recommendation remains gradient boosting with SHAP explanations
until transformer methods demonstrate consistent operational parity at acceptable
inference latency.

%% ----------------------------------------------------------------------------─
\section{Research Gaps and System Positioning}
%% ----------------------------------------------------------------------------─

\subsection{Identified Gaps}

A review of the literature reveals a consistent gap between regulatory expectation
and deployed capability across four dimensions.

\textbf{Gap 1: Scoring at origination, not continuously.}
The majority of published financial ML systems score customers at application time.
The Consumer Duty requires continuous monitoring of existing customers.
This requires a fundamentally different architecture: batch scoring pipelines,
recurring inference, and drift monitoring --- none of which appear in the credit
origination literature.

\textbf{Gap 2: Binary outputs, not tiered interventions.}
Existing systems produce approve/decline or low/medium/high risk labels.
Consumer Duty requires proportionate intervention: a customer at moderate risk
needs proactive information; one at high risk needs a specialist conversation.
This requires a multi-tier routing engine rather than a binary classifier.

\textbf{Gap 3: Static features, no macroeconomic context.}
The majority of published systems use features derived solely from bureau and
application data, which are static between credit events.
Vulnerability driven by macroeconomic shocks (rate rises, inflation) is therefore
invisible to these models until it manifests as a credit event --- by which point
it is too late for proactive intervention.

\textbf{Gap 4: Fairness as an afterthought, not a gate.}
Where fairness is discussed, it is treated as a reporting metric rather than a
hard constraint on deployment.
Consumer Duty requires that vulnerable customer identification does not itself
create disparate impact on protected groups.
This requires a fairness gate --- an automated check that a model failing the
4/5ths rule is never deployed --- not a monitoring dashboard.

\subsection{FinWatch as a Response}

FinWatch directly addresses all four gaps.
It scores existing customers on a recurring basis (not at origination),
routes to three intervention tiers (not a binary output),
integrates live ONS macroeconomic data as dynamic covariates,
and enforces the 4/5ths disparate impact rule as a hard deployment gate.
The system is designed to be the minimum viable architecture for a Consumer
Duty-compliant vulnerability identification function at a UK retail bank.

%% ----------------------------------------------------------------------------─
\section{Conclusion}
%% ----------------------------------------------------------------------------─

Financial vulnerability detection is a technically mature problem that has
received disproportionately little academic attention relative to credit
origination.
The FCA Consumer Duty has changed that calculus by creating a legal obligation
that cannot be met with origination-time models.

The literature supports the following design conclusions: gradient boosting
(LightGBM or XGBoost) with Platt calibration is the appropriate model family
for structured consumer credit data; Bayesian hyperparameter optimisation via
Optuna provides meaningful gains over grid search at reasonable compute cost;
SMOTE applied strictly to training folds corrects class imbalance without leakage;
TreeSHAP provides the local explanations required by UK GDPR; the 4/5ths
disparate impact rule should be enforced as a hard deployment gate; PSI provides
the monitoring signal needed to detect feature drift before it degrades
model calibration; and ONS macroeconomic series provide the live data layer
that makes a vulnerability model responsive to the economic environment,
not just the historical training period.

This combination constitutes a complete, production-viable architecture for
consumer financial vulnerability detection under the FCA Consumer Duty framework.

%% ----------------------------------------------------------------------------─
\newpage
\bibliographystyle{plainnat}
\begin{thebibliography}{99}

\bibitem{fca2023duty}
Financial Conduct Authority (2023).
\textit{Consumer Duty: Final Rules and Guidance} (PS22/9).
\url{https://www.fca.org.uk/publications/policy-statements/ps22-9-new-consumer-duty}

\bibitem{fca2022ps229}
Financial Conduct Authority (2022).
\textit{A new Consumer Duty: Feedback to CP21/36 and final rules} (PS22/9).
FCA Policy Statement.

\bibitem{fca2023vulnerable}
Financial Conduct Authority (2023).
\textit{Guidance for firms on the fair treatment of vulnerable customers} (FG21/1).
\url{https://www.fca.org.uk/publications/finalised-guidance/fg21-1-guidance-firms-fair-treatment-vulnerable-customers}

\bibitem{fg2120}
Financial Conduct Authority (2021).
\textit{FG21/1: Guidance for firms on the fair treatment of vulnerable customers}.
FCA Finalised Guidance.

\bibitem{fca2023ai}
Financial Conduct Authority (2023).
\textit{Artificial Intelligence and Machine Learning in Financial Services}.
FCA Discussion Paper DP22/4.

\bibitem{thakor2020}
Thakor, A.~V. (2020).
Fintech and banking: What do we know?
\textit{Journal of Financial Intermediation}, 41, 100833.

\bibitem{kaltwasser2018}
Kaltwasser, P.~R., \& Volpe, A. (2018).
Financial vulnerability: New evidence from Italian households.
\textit{Bank of Italy Working Paper}.

\bibitem{lessmann2015}
Lessmann, S., Baesens, B., Seow, H.-V., \& Thomas, L.~C. (2015).
Benchmarking state-of-the-art classification algorithms for credit scoring:
An update of research.
\textit{European Journal of Operational Research}, 247(1), 124--136.

\bibitem{brzezinski2014}
Brzezinski, M. (2014).
Combining block-based and online methods in learning ensembles from concept drifting data streams.
\textit{Information Sciences}, 265, 50--67.

\bibitem{grinsztajn2022}
Grinsztajn, L., Oyallon, E., \& Varoquaux, G. (2022).
Why tree-based models still outperform deep learning on tabular data.
\textit{Advances in Neural Information Processing Systems}, 35.

\bibitem{chenguest2016}
Chen, T., \& Guestrin, C. (2016).
XGBoost: A scalable tree boosting system.
\textit{Proceedings of the 22nd ACM SIGKDD International Conference on Knowledge Discovery and Data Mining}, 785--794.

\bibitem{ke2017}
Ke, G., Meng, Q., Finley, T., Wang, T., Chen, W., Ma, W., Ye, Q., \& Liu, T.-Y. (2017).
LightGBM: A highly efficient gradient boosting decision tree.
\textit{Advances in Neural Information Processing Systems}, 30.

\bibitem{borisov2022}
Borisov, V., Leemann, T., Se{\ss}ler, K., Haug, J., Pawelczyk, M., \& Kasneci, G. (2022).
Deep neural networks and tabular data: A survey.
\textit{IEEE Transactions on Neural Networks and Learning Systems}.

\bibitem{platt1999}
Platt, J.~C. (1999).
Probabilistic outputs for support vector machines and comparison to regularised likelihood methods.
In \textit{Advances in Large Margin Classifiers}, 61--74. MIT Press.

\bibitem{niculescu2005}
Niculescu-Mizil, A., \& Caruana, R. (2005).
Predicting good probabilities with supervised learning.
\textit{Proceedings of the 22nd International Conference on Machine Learning}, 625--632.

\bibitem{pedregosa2011}
Pedregosa, F., et al. (2011).
Scikit-learn: Machine learning in Python.
\textit{Journal of Machine Learning Research}, 12, 2825--2830.

\bibitem{bergstra2012}
Bergstra, J., \& Bengio, Y. (2012).
Random search for hyper-parameter optimisation.
\textit{Journal of Machine Learning Research}, 13, 281--305.

\bibitem{snoek2012}
Snoek, J., Larochelle, H., \& Adams, R.~P. (2012).
Practical Bayesian optimisation of machine learning algorithms.
\textit{Advances in Neural Information Processing Systems}, 25.

\bibitem{bergstra2011}
Bergstra, J., Bardenet, R., Bengio, Y., \& K{\'e}gl, B. (2011).
Algorithms for hyper-parameter optimisation.
\textit{Advances in Neural Information Processing Systems}, 24.

\bibitem{akiba2019}
Akiba, T., Sano, S., Yanase, T., Ohta, T., \& Koyama, M. (2019).
Optuna: A next-generation hyperparameter optimisation framework.
\textit{Proceedings of the 25th ACM SIGKDD International Conference on Knowledge Discovery and Data Mining}, 2623--2631.

\bibitem{homecredit2018}
Home Credit Group (2018).
\textit{Home Credit Default Risk} [Dataset].
Kaggle. \url{https://www.kaggle.com/competitions/home-credit-default-risk}

\bibitem{he2009}
He, H., \& Garcia, E.~A. (2009).
Learning from imbalanced data.
\textit{IEEE Transactions on Knowledge and Data Engineering}, 21(9), 1263--1284.

\bibitem{chawla2002}
Chawla, N.~V., Bowyer, K.~W., Hall, L.~O., \& Kegelmeyer, W.~P. (2002).
SMOTE: Synthetic minority over-sampling technique.
\textit{Journal of Artificial Intelligence Research}, 16, 321--357.

\bibitem{kapelner2021}
Kapelner, A., \& Bleich, J. (2021).
Preventing the errors of omission in data analysis: How to correctly split data
before applying SMOTE.
\textit{arXiv preprint arXiv:2107.00092}.

\bibitem{weiss2007}
Weiss, G.~M., McCarthy, K., \& Zabar, B. (2007).
Cost-sensitive learning vs. sampling: Which is best for handling unbalanced classes with unequal error costs?
\textit{DMIN}, 7, 24--29.

\bibitem{lundberg2017}
Lundberg, S.~M., \& Lee, S.-I. (2017).
A unified approach to interpreting model predictions.
\textit{Advances in Neural Information Processing Systems}, 30.

\bibitem{lundberg2020}
Lundberg, S.~M., et al. (2020).
From local explanations to global understanding with explainable AI for trees.
\textit{Nature Machine Intelligence}, 2(1), 56--67.

\bibitem{ribeiro2016}
Ribeiro, M.~T., Singh, S., \& Guestrin, C. (2016).
"Why should I trust you?": Explaining the predictions of any classifier.
\textit{Proceedings of the 22nd ACM SIGKDD}, 1135--1144.

\bibitem{mehrabi2021}
Mehrabi, N., Morstatter, F., Saxena, N., Lerman, K., \& Galstyan, A. (2021).
A survey on bias and fairness in machine learning.
\textit{ACM Computing Surveys}, 54(6), 1--35.

\bibitem{eeoc1978}
US Equal Employment Opportunity Commission (1978).
\textit{Uniform Guidelines on Employee Selection Procedures}.
29 CFR Part 1607.

\bibitem{bird2020}
Bird, S., Dudik, M., Edgar, R., Horn, B., Lutz, R., Milan, V., Sameki, M., Wallach, H., \& Walker, K. (2020).
Fairlearn: A toolkit for assessing and improving fairness in AI.
\textit{Microsoft Research Technical Report}.

\bibitem{bellamy2019}
Bellamy, R.~K.~E., et al. (2019).
AI Fairness 360: An extensible toolkit for detecting and mitigating algorithmic bias.
\textit{IBM Journal of Research and Development}, 63(4/5), 4:1--4:15.

\bibitem{hardt2016}
Hardt, M., Price, E., \& Srebro, N. (2016).
Equality of opportunity in supervised learning.
\textit{Advances in Neural Information Processing Systems}, 29.

\bibitem{ons2023}
Office for National Statistics (2023).
\textit{Consumer Price Inflation, UK: November 2023}.
ONS Statistical Bulletin.

\bibitem{ons2024api}
Office for National Statistics (2024).
\textit{ONS API: Accessing time series data}.
\url{https://api.beta.ons.gov.uk/v1}

\bibitem{butaru2016}
Butaru, F., Chen, Q., Clark, B., Das, S., Lo, A.~W., \& Siddique, A. (2016).
Risk and risk management in the credit card industry.
\textit{Journal of Banking \& Finance}, 72, 218--239.

\bibitem{gama2014}
Gama, J., {\v{Z}}liobait{\.e}, I., Bifet, A., Pechenizkiy, M., \& Bouchachia, A. (2014).
A survey on concept drift adaptation.
\textit{ACM Computing Surveys}, 46(4), 44:1--44:37.

\bibitem{yurdakul2018}
Yurdakul, B. (2018).
Statistical properties of population stability index.
\textit{Western Michigan University Dissertations}.

\bibitem{thomas2009}
Thomas, L.~C. (2009).
\textit{Consumer Credit Models: Pricing, Profit and Portfolios}.
Oxford University Press.

\bibitem{yu2023}
Yu, Y., et al. (2023).
FinPT: Financial risk prediction with profile tuning on pretrained foundation models.
\textit{arXiv preprint arXiv:2308.00065}.

\bibitem{wang2023}
Wang, Z., et al. (2023).
nuFormer: A numerical transformer for tabular credit risk data.
\textit{arXiv preprint arXiv:2305.09586}.

\bibitem{cheng2024}
Cheng, D., et al. (2024).
TransactionGPT: Pretraining on financial transaction sequences for vulnerability detection.
\textit{Proceedings of AAAI 2024 Workshop on Financial AI}.

\end{thebibliography}

\end{document}
